% =====================================================================
%  Shared preamble of the AB02 documents (complete documentation and
%  mathematical models). \input from both main files.
% =====================================================================
\IfFileExists{lmodern.sty}{\usepackage[T1]{fontenc}\usepackage{lmodern}}{}
\usepackage{microtype}
\usepackage{amsmath,amssymb,amsthm,bm,mathtools}
\usepackage{graphicx}
\usepackage{booktabs,array,tabularx,longtable,multirow}
\usepackage{xcolor}
\usepackage[most]{tcolorbox}
\usepackage{enumitem}
\usepackage{fancyhdr}
\usepackage{titlesec}
\usepackage{caption}
\usepackage{float}
\usepackage{listings}
\usepackage{needspace}
\usepackage{tikz}
\usetikzlibrary{arrows.meta, positioning, calc, fit, backgrounds}
\usepackage[hidelinks]{hyperref}
\usepackage[nameinlink]{cleveref}

% ---------------------------------------------------------------- colours
\definecolor{cHead}{HTML}{0D3B66}
\definecolor{cText}{HTML}{263238}
\definecolor{cPlant}{HTML}{37474F}
\definecolor{cEnv}{HTML}{A0620F}
\definecolor{cCtrl}{HTML}{AD1457}
\definecolor{cAdapt}{HTML}{4527A0}
\definecolor{cX3}{HTML}{1565C0}
\definecolor{cOk}{HTML}{2E7D32}
\definecolor{cWarn}{HTML}{B23A48}
\definecolor{cCodeBg}{HTML}{F4F4F4}
\definecolor{cKey}{HTML}{0000FF}
\definecolor{cCom}{HTML}{228B22}
\definecolor{cStr}{HTML}{A020F0}
\hypersetup{colorlinks=true, linkcolor=cHead, urlcolor=cX3, citecolor=cHead}

% ---------------------------------------------------------------- headings
\titleformat{\section}{\Large\sffamily\bfseries\color{cHead}}{\thesection}{0.8em}{}[{\color{cHead}\titlerule[0.8pt]}]
\titleformat{\subsection}{\large\sffamily\bfseries\color{cHead}}{\thesubsection}{0.7em}{}
\titleformat{\subsubsection}{\normalsize\sffamily\bfseries\color{cText}}{\thesubsubsection}{0.6em}{}
\titlespacing*{\section}{0pt}{2.2ex plus 1ex}{1.4ex}
\titlespacing*{\subsection}{0pt}{2.0ex plus 1ex}{1.0ex}

\setlist{itemsep=2pt, topsep=4pt}
\setlength{\parindent}{0pt}
\setlength{\parskip}{5pt}
\captionsetup{font=small, labelfont={bf,sf}}
\renewcommand{\arraystretch}{1.15}

% ---------------------------------------------------------------- notation
\newcommand{\vq}{\bm{q}}  \newcommand{\qv}{\bm{q}_v} \newcommand{\vw}{\bm{\omega}}
\newcommand{\vu}{\bm{u}}  \newcommand{\vd}{\bm{d}}   \newcommand{\vs}{\bm{s}}
\newcommand{\vtau}{\bm{\tau}} \newcommand{\veps}{\bm{\varepsilon}} \newcommand{\vF}{\bm{F}}
\newcommand{\vp}{\bm{p}}  \newcommand{\vv}{\bm{v}}  \newcommand{\ve}{\bm{e}}
\newcommand{\vx}{\bm{x}}  \newcommand{\vr}{\bm{r}}
\newcommand{\that}{\hat{\theta}} \newcommand{\ttil}{\tilde{\theta}}
\newcommand{\skw}[1]{#1^{\times}}
\newcommand{\R}{\mathbb{R}}
\newcommand{\sat}{\operatorname{sat}}
\newcommand{\sgn}{\operatorname{sgn}}
\newcommand{\diag}{\operatorname{diag}}
\newcommand{\umax}{u_{\max}}
\newcommand{\cbreakus}{\textunderscore\allowbreak}
\newcommand{\code}[1]{{\ttfamily\small\let\_\cbreakus #1}}

% ---------------------------------------------------------------- boxes
\tcbset{base/.style={enhanced, breakable, boxrule=0.7pt, arc=2.5pt, left=7pt, right=7pt,
  top=5pt, bottom=5pt, fonttitle=\sffamily\bfseries\small, before skip=8pt, after skip=8pt,
  before upper={\setlength{\parskip}{4pt}}}}
\newtcolorbox{eqbox}[2]{base, colback=#1!4, colframe=#1, colbacktitle=#1, coltitle=white, title={#2}}
\newtcolorbox{toybox}[1][]{base, colback=cEnv!5, colframe=cEnv!70, colbacktitle=cEnv!14,
  coltitle=cEnv!80!black, title={Toy example\ifx&#1&\else\ --- #1\fi}}
\newtcolorbox{idea}[1][]{base, colback=cX3!4, colframe=cX3!55, colbacktitle=cX3!12,
  coltitle=cX3!80!black, title={\ifx&#1&The idea in plain words\else #1\fi}}
\newtcolorbox{projbox}[1][]{base, colback=cOk!5, colframe=cOk!70, colbacktitle=cOk!14,
  coltitle=cOk!70!black, title={In our project\ifx&#1&\else\ --- #1\fi}}
\newtcolorbox{warn}[1][]{base, colback=cWarn!5, colframe=cWarn!80, colbacktitle=cWarn!15,
  coltitle=cWarn!80!black, title={\ifx&#1&Be careful\else #1\fi}}
\newtcolorbox{codebox}{base, colback=cPlant!4, colframe=cPlant!45, boxrule=0.5pt,
  fontupper=\small, left=6pt, top=3pt, bottom=3pt}

% symbol table:  \begin{dimtable} sym & meaning & dimension & unit \\ \end{dimtable}
\newenvironment{dimtable}{%
  \par\smallskip\small
  \tabularx{\linewidth}{@{}>{$}l<{$} X >{$}l<{$} l@{}}
  \toprule
  \textrm{\textbf{Symbol}} & \textbf{Meaning} & \textrm{\textbf{Dimension}} & \textbf{Unit}\\
  \midrule}{%
  \bottomrule\endtabularx\par\smallskip}

% ---------------------------------------------------------------- code listings
\lstdefinestyle{common}{basicstyle=\ttfamily\footnotesize, backgroundcolor=\color{cCodeBg},
  frame=single, rulecolor=\color{black!20}, framesep=4pt, xleftmargin=6pt, xrightmargin=2pt,
  keywordstyle=\color{cKey}, commentstyle=\color{cCom}, stringstyle=\color{cStr},
  showstringspaces=false, breaklines=true, columns=fullflexible, keepspaces=true,
  numberstyle=\tiny\color{black!50}, aboveskip=6pt, belowskip=6pt, upquote=true}
\lstdefinestyle{py}{style=common, language=Python, numbers=left, numbersep=6pt}
\lstdefinestyle{mat}{style=common, language=Matlab}
\lstdefinestyle{sh}{style=common, language=bash}

% ---------------------------------------------------------------- toy-model figure
% \toyfig{file.png}{width}{caption}{label}: figure produced by toy_model/run_all.m (MATLAB).
\newcommand{\toyfig}[4]{%
  \begin{figure}[H]\centering
  \IfFileExists{../../toy_model/figures/#1}{\includegraphics[width=#2]{../../toy_model/figures/#1}}{%
    \fbox{\parbox{0.8\linewidth}{\centering\vspace{1.2cm}\sffamily\small
      \texttt{toy\_model/figures/\detokenize{#1}}\\ is created by \texttt{toy\_model/run\_all.m} (MATLAB).\vspace{1.2cm}}}}
  \caption{#3 \emph{(MATLAB toy model, \texttt{toy\_model/figures/\detokenize{#1}}.)}}\label{#4}
  \end{figure}}
% \resfig{file.png}{width}{caption}{label}: figure produced by the Python scripts (docs/figures).
\newcommand{\resfig}[4]{%
  \begin{figure}[H]\centering\includegraphics[width=#2]{../figures/#1}
  \caption{#3}\label{#4}\end{figure}}

% ---------------------------------------------------------------- code read straight from the repository
% \pyfile{first line}{last line}{file}{title}
\newcommand{\pyfile}[4]{\lstinputlisting[style=py, firstline=#1, lastline=#2, firstnumber=#1,
  title={\small\sffamily #4}]{#3}}
\newcommand{\matfile}[4]{\lstinputlisting[style=mat, numbers=left, numbersep=6pt, firstline=#1, lastline=#2,
  firstnumber=#1, title={\small\sffamily #4}]{#3}}
\newcommand{\repro}{../../paper_reproduction.py}
\newcommand{\xmodel}{../../x3_sim/x3_model.py}
\newcommand{\xctrl}{../../x3_sim/controllers.py}
\newcommand{\xrun}{../../x3_sim/run_x3_experiments.py}
\newcommand{\gznode}{../../gazebo.py}
\newcommand{\toylib}{../../toy_model/lib}
\newcommand{\toydir}{../../toy_model}
